\chapter{Government Debt and Ricardian Equivalence}\label{ch:b2:government-debt}

Governments borrow, sometimes a great deal. The chart shows the debt of five governments as a share of their country's
GDP over a century and a half.

\begin{openingchart}{General government debt as a share of GDP, per cent, United Kingdom, United States, France, Japan
  and Italy, 1870--2020; gaps where the series are missing.}{Jord\`a--Schularick--Taylor Macrohistory Database (debtgdp)}
\begin{tikzpicture}
\begin{axis}[omrecord, width=0.86\linewidth, height=0.44\linewidth, xmin=1870, xmax=2020, ymin=0, ymax=280,
    xlabel={Year}, ylabel={Per cent of GDP}, xtick={1880,1900,1920,1940,1960,1980,2000,2020},
    xticklabel style={/pgf/number format/1000 sep=}, unbounded coords=jump, legend style={at={(0.5,-0.3)},
    anchor=north, legend columns=5, /tikz/every even column/.append style={column sep=8pt}}]
  \addplot[omDef, thick] table[x=year, y=GBR] {figdata/bachelor-2/21-government-debt/out/debt.csv};
  \addplot[omThm, thick, dashed] table[x=year, y=USA] {figdata/bachelor-2/21-government-debt/out/debt.csv};
  \addplot[omMeth, thick, dotted] table[x=year, y=FRA] {figdata/bachelor-2/21-government-debt/out/debt.csv};
  \addplot[black!70, thick, dashdotted] table[x=year, y=JPN] {figdata/bachelor-2/21-government-debt/out/debt.csv};
  \addplot[omProp, thick] table[x=year, y=ITA] {figdata/bachelor-2/21-government-debt/out/debt.csv};
  \legend{United Kingdom, United States, France, Japan, Italy}
\end{axis}
\end{tikzpicture}
\end{openingchart}

% figdata: debt_summary (GBR peak 270 in 1946; FRA 237 in 1921; 2020: GBR 104, USA 128, FRA 115, JPN 254, ITA 156)
The peaks before 1950 come from wars: British debt reached 270 per cent of GDP in 1946, French debt 237 per cent in 1921. After each
war the ratios fell, sometimes fast. The peaks of 2020 are of another kind: Japan's debt was 254 per cent of GDP, Italy's 156,
the United States' 128, all reached in peacetime. This chapter asks how debt evolves, when it can be sustained, whether
borrowing instead of taxing changes what households do, and how large debts have come down.

\begin{recall}
From the school volumes: the government budget, the deficit and the debt. From the Year 1 volume: crowding out in the
long-run model. From \cref{ch:b2:consumption}: the permanent-income model and \omterm{def:b2:consumption:hand-to-mouth}{hand-to-mouth households}.
\end{recall}

\begin{notation}[Symbols of this chapter]\label{not:b2:government-debt:symbols}
$B_t$ is government debt at the end of year $t$ and $b_t = B_t/Y_t$ its ratio to GDP; $S_t$ the primary surplus and $s_t = S_t/Y_t$; $i$ the nominal interest rate paid on
the debt; $\gamma$ the growth rate of nominal GDP; $T_t$ taxes and $G_t$ government purchases.
\end{notation}

\section{The budget constraint over time}\label{sec:b2:government-debt:budget}

\begin{definition}[Primary balance]\label{def:b2:government-debt:primary-balance}
The \emph{primary balance}\index{primary balance} is government revenue minus government spending other than interest on
the debt. A positive primary balance is a primary surplus.
\end{definition}

Each year the debt grows by the interest paid on it and falls by the primary surplus: $B_{t+1} = (1 + i)B_t - S_{t+1}$. Dividing by $Y_{t+1} = (1 + \gamma)Y_t$ gives the
law of motion of the debt ratio:
\begin{equation}\label{eq:b2:government-debt:dynamics}
b_{t+1} = \frac{1 + i}{1 + \gamma}\,b_t - s_{t+1}.
\end{equation}
Two forces move the ratio: the gap between the interest rate and the growth rate of nominal GDP, which acts on the whole
stock of debt, and the primary surplus.

\begin{definition}[Intertemporal government budget constraint]\label{def:b2:government-debt:igbc}
The \emph{intertemporal government budget constraint}\index{intertemporal government budget constraint} states that the
debt today equals the present value of future primary surpluses, discounted at the interest rate on the debt: when
$i > \gamma$, $b_0 = \sum_{t \ge 1} s_t\left(\frac{1 + \gamma}{1 + i}\right)^t$.
\end{definition}

It follows from \eqref{eq:b2:government-debt:dynamics} by solving forward, provided the discounted debt $b_T((1 + \gamma)/(1 + i))^T$ tends to zero: lenders will not let a government repay
its debt for ever with new debt. A government that owes debt must, in present value, run primary surpluses. When $i < \gamma$ the
discounted debt shrinks without surpluses, and the constraint loses its force; \cref{ch:b2:overlapping-generations} returns to this case.

\section{Sustainability}\label{sec:b2:government-debt:sustainability}

\begin{definition}[Debt-stabilising primary surplus]\label{def:b2:government-debt:stabilising-surplus}
The \emph{debt-stabilising primary surplus}\index{debt-stabilising primary surplus} is the primary surplus, as a share
of GDP, that keeps the debt ratio constant: $s^* = \frac{i - \gamma}{1 + \gamma}\,b$.
\end{definition}

With debt at 90 per cent of GDP, an interest rate of 4 per cent and nominal growth of 3 per cent, $s^* = 0.01/1.03 \times 0.9 = 0.874$ per cent of GDP. If growth exceeded the
interest rate, the stabilising balance would be a primary deficit.

\begin{omfigure}
\begin{tikzpicture}[econ]
  \begin{scope}
    \draw[->] (0,-1.4) -- (0,1.6) node[left] {$b_{t+1} - b_t$};
    \draw[->] (0,0) -- (4.4,0) node[below] {$b_t$};
    \draw[omDef, thick] (0,-1.2) -- (4,1.3);
    \node[eqpoint] at (1.92,0) {};
    \node[below right, font=\footnotesize] at (1.92,0) {$b^*$};
    \draw[->, thick] (1.6,0.12) -- (0.9,0.12);
    \draw[->, thick] (2.3,0.12) -- (3.0,0.12);
    \node[font=\footnotesize] at (2,-1.9) {$i > \gamma$: unstable};
  \end{scope}
  \begin{scope}[xshift=6cm]
    \draw[->] (0,-1.4) -- (0,1.6) node[left] {$b_{t+1} - b_t$};
    \draw[->] (0,0) -- (4.4,0) node[below] {$b_t$};
    \draw[omThm, thick] (0,1.2) -- (4,-1.3);
    \node[eqpoint] at (1.92,0) {};
    \node[below left, font=\footnotesize] at (1.92,0) {$b^*$};
    \draw[->, thick] (0.9,0.12) -- (1.6,0.12);
    \draw[->, thick] (3.0,0.12) -- (2.3,0.12);
    \node[font=\footnotesize] at (2,-1.9) {$i < \gamma$: stable};
  \end{scope}
\end{tikzpicture}

{\small Debt dynamics with a constant \omterm{def:b2:government-debt:primary-balance}{primary balance}. The change in the ratio is $\frac{i - \gamma}{1 + \gamma}b_t - s$. When $i > \gamma$ (left) the line slopes up: above the ratio $b^*$ that
the surplus stabilises, debt grows without limit unless the surplus rises. When $i < \gamma$ (right) the line slopes down: the ratio
converges to $b^*$ from any start, even with a primary deficit.}
\end{omfigure}

\begin{definition}[Fiscal sustainability]\label{def:b2:government-debt:sustainability}
Fiscal policy is \emph{sustainable}\index{fiscal sustainability} if, continued, it satisfies the intertemporal budget
constraint without a default or an unplanned change of policy.
\end{definition}

\omterm{def:b2:government-debt:sustainability}{Sustainability} is a property of policy, not of a debt ratio, and it cannot be read from one year's numbers. A practical
test asks whether governments respond to debt: if the primary surplus rises when the debt ratio rises, by more than $(i - \gamma)/(1 + \gamma)$
for each point of debt, the ratio is pulled back.

% ledger: b2:bohn1998
Bohn (1998) ran this test on the United States for 1916--1995. Unit-root tests on the debt ratio alone could not reject that
it wanders without returning; but the primary surplus rose with the debt ratio, with a coefficient of 0.054 over the full
sample: each 100 dollars of extra debt raised the next year's primary surplus by 5.40 dollars. Across his samples the
coefficient lay between 2.8 and 5.4 per cent.

The listing computes, from Eurostat's government accounts, the stabilising \omterm{def:b2:government-debt:primary-balance}{primary balance} of four European governments
in 2024 and compares it with the actual one.

\omcode{code/bachelor-2/21-government-debt/listings/stabilise.py}{1}{30}{The debt-stabilising \omterm{def:b2:government-debt:primary-balance}{primary balance} and the
actual \omterm{def:b2:government-debt:primary-balance}{primary balance} of four governments, 2024.}\label{lst:b2:government-debt:stabilise}

\omoutput{code/bachelor-2/21-government-debt/outputs/stabilise.txt}{In 2024 Germany and France ran \omterm{def:b2:government-debt:primary-balance}{primary balances}
below the stabilising one; Italy and Greece ran primary surpluses well above it.}\label{out:b2:government-debt:stabilise}

% code: listings/stabilise.py
The table rests on an average: from 2019 to 2024 nominal GDP grew fast, in part because prices rose fast, while the
interest paid on the stock of debt, fixed for the life of each bond, adjusted only as old bonds were replaced. Over that window growth
exceeded the interest rate everywhere, so every government could have run a primary deficit and held its ratio. With
another window, or when old debt is refinanced at higher rates, the answer changes; an exercise tries it.

\section{Ricardian equivalence}\label{sec:b2:government-debt:ricardian}

Suppose the government cuts taxes today and borrows to pay for the cut, planning to raise taxes later to repay. Does
consumption rise?

\begin{definition}[Ricardian equivalence]\label{def:b2:government-debt:ricardian-equivalence}
\emph{Ricardian equivalence}\index{Ricardian equivalence} is the proposition that, for a given path of government
purchases, replacing taxes by debt, or debt by taxes, leaves households' consumption unchanged: households save the tax
cut to pay the future taxes it implies.
\end{definition}

\begin{proposition}[Ricardian equivalence in two periods]\label{prop:b2:government-debt:ricardian}
Let a household live two periods with incomes $y_1$, $y_2$ and pay taxes $T_1$, $T_2$, and let it borrow and lend at the rate $r$ at which the government
borrows. If the government cuts $T_1$ by $\Delta$ and raises $T_2$ by $(1 + r)\Delta$ to repay the debt, the household's consumption in both periods is unchanged and its saving rises by $\Delta$.
\end{proposition}

\begin{proof}
The household's budget constraint is $c_1 + c_2/(1 + r) = y_1 - T_1 + (y_2 - T_2)/(1 + r)$. The change in the right-hand side is $\Delta - (1 + r)\Delta/(1 + r) = 0$. The set of affordable plans is
unchanged, so the chosen plan is unchanged. With $c_1$ the same and disposable income $y_1 - T_1$ higher by $\Delta$, saving $y_1 - T_1 - c_1$ rises by $\Delta$.
\end{proof}

\begin{omfigure}
\begin{tikzpicture}[econ]
\begin{axis}[omrecord, grid=none, xmin=0, xmax=210, ymin=0, ymax=220, width=0.52\linewidth, height=0.46\linewidth,
    xlabel={$c_1$}, ylabel={$c_2$}, xtick=\empty, ytick=\empty]
  \addplot[black, thick, domain=0:200, samples=2] {210 - 1.05*x};
  \addplot[omDef, thick, domain=55:190, samples=60] {10500/x};
  \node[eqpoint] at (axis cs:100,105) {};
  \node[anchor=south west, font=\footnotesize] at (axis cs:100,105) {$C$};
  \node[eqpoint, fill=omThm] at (axis cs:120,84) {};
  \node[anchor=north east, font=\footnotesize] at (axis cs:119,83) {$E$};
  \node[eqpoint, fill=omThm] at (axis cs:140,63) {};
  \node[anchor=north east, font=\footnotesize] at (axis cs:139,62) {$E'$};
\end{axis}
\end{tikzpicture}

{\small A debt-financed tax cut in two periods. Disposable income moves from $E$ to $E'$: more now, less later, with the same present value,
so along the same budget line. The household's choice $C$ does not move; it saves the extra disposable income.}
\end{omfigure}

The proof uses four assumptions, and each is a reason the proposition may fail. Households must be able to borrow and lend
at the government's rate: a household that cannot borrow, the \omterm{def:b2:consumption:hand-to-mouth}{hand-to-mouth household} of \cref{ch:b2:consumption}, spends a tax cut. The
future taxes must fall on the same households: if they fall on people not yet born, and the living do not care for them,
a tax cut makes the living richer. Taxes must be lump-sum: a tax on income changes incentives, so moving it in time changes
behaviour. And households must foresee the future taxes.

\begin{model}{Ricardian equivalence}
\assumes Households that borrow and lend at the government's rate, care about the households that will pay future taxes,
foresee those taxes, and pay lump-sum taxes; a given path of government purchases.
\predicts A debt-financed tax cut raises private saving one for one, and leaves national saving and consumption unchanged.
\tested Across countries, changes in government net lending are partly offset in national saving: about half of a change
passes into national saving (the chart below).
\breaks With borrowing limits and finite horizons: households spend part of a tax cut, all of it for \omterm{def:b2:consumption:hand-to-mouth}{hand-to-mouth
households}.
\end{model}

If \omterm{def:b2:government-debt:ricardian-equivalence}{Ricardian equivalence} held, a rise in the government deficit would be matched by a rise in private saving, and national
saving would not change. If households ignored future taxes entirely, national saving would fall one for one with
government saving. The chart compares, across 87 countries, the change in national saving with the change in government net
lending between the years before and after the financial crisis.

\begin{omchart}{Change in gross national saving (per cent of GNI) against change in general government net lending (per
  cent of GDP), average 2010--2017 minus average 2000--2007, 87 countries; least-squares line, and the line of no offset
  (slope 1).}{World Bank, World Development Indicators (NY.GNS.ICTR.ZS, GC.NLD.TOTL.GD.ZS); our computation}
\begin{tikzpicture}
\begin{axis}[omrecord, width=0.66\linewidth, height=0.5\linewidth, xmin=-14, xmax=14, ymin=-15, ymax=26,
    xlabel={Change in government net lending, points}, ylabel={Change in national saving, points},
    legend style={at={(0.5,-0.25)}, anchor=north, legend columns=2, /tikz/every even column/.append style={column
    sep=10pt}}]
  \addplot[only marks, mark=*, mark size=1.3pt, omDef, fill opacity=0.6, forget plot]
    table[x=dn, y=ds] {figdata/bachelor-2/21-government-debt/out/offset.csv};
  \addplot[omThm, thick, domain=-14:14, samples=2] {0.48 + 0.53*x};
  \addplot[black!60, dashed, domain=-14:14, samples=2] {0.48 + x};
  \legend{Fitted (slope 0.53), No offset (slope 1)}
\end{axis}
\end{tikzpicture}
\end{omchart}

% figdata: offset_summary (n 87, slope 0.53, se 0.19, intercept 0.48, corr 0.3)
The slope is 0.53, with a standard error of 0.19: about half of a change in government saving passed into national saving, and
about half was offset by private saving. Both full offset (slope 0) and no offset (slope 1) are rejected. The test is rough: government
net lending includes public investment, which Ricardian reasoning treats differently, and the decade also saw changes in
private saving with other causes, which is why the correlation is only 0.3.

\section{Tax smoothing}\label{sec:b2:government-debt:tax-smoothing}

\begin{definition}[Tax smoothing]\label{def:b2:government-debt:tax-smoothing}
\emph{Tax smoothing}\index{tax smoothing} is the policy of keeping tax rates roughly constant over time, using debt to
absorb temporary changes in spending or revenue, because the efficiency cost of a tax rises more than in proportion to
its rate.
\end{definition}

If the deadweight loss of a tax grows with the square of its rate, as the deadweight-loss triangle suggests, then two years of
a 10 per cent tax cost $10^2 + 10^2 = 200$ units of loss against $0^2 + 20^2 = 400$ for one year at 0 and one at 20, though both raise the same revenue (with a zero interest rate). A government
facing a war, which raises spending for a few years, should then borrow for the war and raise taxes a little for a long
time to repay.

\begin{model}{Tax smoothing}
\assumes A government that minimises the efficiency cost of the taxes needed to finance a given path of spending, with
convex costs of taxation and access to borrowing.
\predicts Debt rises in wars and other temporary spending surges and falls slowly afterwards; tax rates move less than
spending.
\tested Only consistency is shown here: debt rose in the two world wars and fell in the decades after them (the opening
chart). A test needs the model's distinctive prediction, tax rates that move less than spending, checked against
series of tax rates and spending.
\breaks When debt rises in peacetime with no temporary surge in spending, as in Japan and Italy since 1980 (the opening
chart), the model does not explain it.
\end{model}

\section{How large debts ended}\label{sec:b2:government-debt:ending}

Equation \eqref{eq:b2:government-debt:dynamics} lists the ways a debt ratio can fall: primary surpluses, growth of real GDP, inflation that raises nominal GDP faster
than the interest rate, a lower interest rate, or a default that reduces $B$ itself.

\begin{definition}[Financial repression]\label{def:b2:government-debt:financial-repression}
\emph{Financial repression}\index{financial repression} is a set of regulations that hold the interest rates a
government pays below the rates savers would otherwise obtain, for example by requiring banks and pension funds to hold
government bonds, capping interest rates or restricting the movement of capital abroad.
\end{definition}

\begin{record}[How Britain's debt fell after 1945]
% figdata: uk_summary (debt 270 -> 47; nominal 8.9, real_total 2.4, deflator 6.3, inflation 5.7, ltrate 6.4)
British government debt fell from 270 per cent of GDP in 1946 to 47 per cent in 1975. Over those years nominal GDP grew by 8.9 per cent a year on
average: real GDP by 2.4 per cent and its prices, the GDP deflator, by 6.3 per cent (consumer prices rose 5.7 per cent a year), while the long-term government bond rate averaged 6.4
per cent. Nominal growth above the bond rate shrank the ratio year after year; how much primary surpluses added cannot be read
from this database, which does not report them.
\end{record}

\begin{record}[Germany's debt agreement of 1953]
% ledger: b2:london1953
The London Debt Agreement of 1953 settled Germany's debts from the period between the two world wars. It cut the overall
debt by roughly half, stretched repayment over a longer period, and let Germany defer some payments until reunification
(Guinnane 2004). A debt reduced by agreement is a default by another name: the creditors accepted less than they had
been owed.
\end{record}

\begin{debate}[Does high public debt slow growth?]
% ledger: b2:rr2010, b2:hap2014, b2:pss2014
Reinhart and Rogoff (2010) found that above a debt of 90 per cent of GDP, median growth rates were about one percentage point
lower, and average growth considerably lower. Herndon, Ash and Pollin (2014) replicated the calculation for 20 advanced
economies after the war and found coding errors, the exclusion of available data and unusual weighting: computed
correctly, average growth above 90 per cent was 2.2 per cent, not $-0.1$, not dramatically different from growth at lower debt.
Pescatori, Sandri and Simon (2014) found no threshold above which medium-term growth is dramatically compromised, and found
that countries with high but declining debt grew about as fast as those with lower debt. Even a negative correlation would
leave the direction open: slow growth raises debt ratios, as well as the reverse. How much, if at all, high debt lowers
growth is not settled.
\end{debate}

\section{Exercises}\label{sec:b2:government-debt:exercises}

\begin{exercise}[$\star$]\label{exo:b2:government-debt:1}
Read the opening chart. When did each of the five ratios reach its peak, and which peaks came in wartime?
\end{exercise}

\begin{exercise}[$\star$]\label{exo:b2:government-debt:2}
A government runs a deficit of 3 per cent of GDP and pays interest worth 2.5 per cent of GDP. What is its \omterm{def:b2:government-debt:primary-balance}{primary balance}?
\end{exercise}

\begin{exercise}[$\star$]\label{exo:b2:government-debt:3}
Read the offset chart. What slope would \omterm{def:b2:government-debt:ricardian-equivalence}{Ricardian equivalence} predict, what slope would households that ignore future
taxes produce, and what does the fitted slope say?
\end{exercise}

\begin{exercise}[$\star$]\label{exo:b2:government-debt:4}
Describe an observation on households that would falsify \omterm{def:b2:government-debt:ricardian-equivalence}{Ricardian equivalence}, and one that would support it.
\end{exercise}

\begin{exercise}[$\star\star$]\label{exo:b2:government-debt:5}
Predict, then compute: with debt at 90 per cent of GDP and growth of 3 per cent, the stabilising primary surplus at an interest rate of 4 per cent and at 3 per cent.
\end{exercise}

\begin{exercise}[$\star\star$]\label{exo:b2:government-debt:6}
With $i = 2$ per cent and $\gamma = 4$ per cent, a government runs a primary deficit of 2 per cent of GDP for ever. Show that the debt ratio converges, and find its limit.
\end{exercise}

\begin{exercise}[$\star\star$]\label{exo:b2:government-debt:7}
\textbf{(D)} Extend the listing to Spain and Portugal. Then recompute the six stabilising balances with nominal growth
averaged over 2014--2024. Which conclusions change?
\end{exercise}

\begin{exercise}[$\star\star$]\label{exo:b2:government-debt:8}
A household with income 100 in each of two periods pays taxes of 20 in each. The government cuts the first tax by 10 and
borrows at 5 per cent. Find the second-period tax, and the change in the household's consumption and saving if it can
borrow at 5 per cent and if it cannot borrow and wants to consume more now.
\end{exercise}

\begin{exercise}[$\star\star$]\label{exo:b2:government-debt:9}
A government must spend 20 in one of two years and nothing in the other; the interest rate is zero, and the deadweight
loss of a tax rate $\tau$ is proportional to $\tau^2$. Compare the losses of taxing in the year of spending and of smoothing.
\end{exercise}

\begin{exercise}[$\star\star\star$]\label{exo:b2:government-debt:10}
Derive \eqref{eq:b2:government-debt:dynamics} from $B_{t+1} = (1 + i)B_t - S_{t+1}$, the stabilising surplus from it, and the intertemporal budget constraint by solving forward when
$i > \gamma$.
\end{exercise}

\begin{exercise}[$\star\star\star$]\label{exo:b2:government-debt:11}
Let the household of the proposition live only one period, and let the debt be repaid by the next generation, whose
welfare it does not weigh. Show that the tax cut raises its consumption, and explain why it would not if it cared for its
heirs as for itself.
\end{exercise}

\begin{exercise}[$\star\star\star$]\label{exo:b2:government-debt:12}
Let the primary surplus respond to debt, $s_{t+1} = \rho b_t + \mu$. Show that the debt ratio converges if $\rho > (i - \gamma)/(1 + \gamma)$. With Bohn's $\rho = 0.054$, $i = 4$ per cent and $\gamma = 3$ per cent, find the
half-life of a deviation of the ratio from its long-run level.
\end{exercise}

\section{Problem: Brenland's debt}\label{sec:b2:government-debt:problem}

\begin{problem}[{Weekend problem --- Brenland's debt: its accounts, the surplus that holds its ratio, a debt-financed tax cut,
and the years to reach a lower ratio}]
\label{pb:b2:government-debt:1}
Use Brenland's public accounts (\texttt{accounts.csv}, millions of dalens). In 2024 the debt was 228,953 and GDP 379,562; in 2025 GDP was 393,264, the deficit 7,865
and the debt 236,818. The interest rate paid on the debt is 3 per cent.

\medskip
\noindent\textbf{Part I --- The accounts.}
\begin{enumerate}
  \item What was the deficit in 2025 as a share of GDP?
  \item What was the debt ratio at the end of 2025?
  \item Check that the debt rose by the deficit.
  \item What interest did the government pay in 2025?
  \item What was its \omterm{def:b2:government-debt:primary-balance}{primary balance}, and as a share of GDP?
  \item In 2022 the deficit was 2 per cent of GDP, but the debt ratio fell from 68.5 to 63.8 per cent. How?
\end{enumerate}

\medskip
\noindent\textbf{Part II --- Dynamics.} From now on nominal GDP grows 3.6 per cent a year; take $b = 0.602$.
\begin{enumerate}[resume]
  \item What is the stabilising \omterm{def:b2:government-debt:primary-balance}{primary balance} at an interest rate of 3 per cent?
  \item With the 2025 primary deficit kept, does the ratio rise or fall?
  \item What is the ratio after one more year?
  \item If the interest rate rose to 5 per cent, what would the stabilising balance be?
  \item And the ratio after one year, with the same primary deficit?
  \item Why does the interest rate act on the whole debt, and the \omterm{def:b2:government-debt:primary-balance}{primary balance} only on the year's flow?
\end{enumerate}

\medskip
\noindent\textbf{Part III --- A tax cut.} The government cuts each household's taxes by 1,000 dalens this year, borrows at 3 per
cent, and raises the tax next year to repay.
\begin{enumerate}[resume]
  \item By how much does next year's tax rise?
  \item What does a household that can borrow and lend at 3 per cent do with the cut?
  \item What does a \omterm{def:b2:consumption:hand-to-mouth}{hand-to-mouth household} do?
  \item If 30 per cent of households are hand to mouth, by how much does private saving rise per household, on average?
  \item By how much does national saving fall, and what share of the cut is offset?
  \item Name two assumptions of \omterm{def:b2:government-debt:ricardian-equivalence}{Ricardian equivalence} that the \omterm{def:b2:consumption:hand-to-mouth}{hand-to-mouth household} violates or might violate.
\end{enumerate}

\medskip
\noindent\textbf{Part IV --- A lower ratio.} Brenland aims at a debt ratio of 50 per cent.
\begin{enumerate}[resume]
  \item With a primary surplus of 1 per cent of GDP, at 3 per cent interest, how many years does it take?
  \item At 5 per cent interest?
  \item What constant primary surplus reaches 50 per cent in ten years, at 3 per cent?
  \item If inflation raised nominal growth to 5.6 per cent while the interest rate stayed at 3 per cent, how many years would it take?
  \item What would the government need to keep the interest rate at 3 per cent in that case, and what would savers lose?
  \item State the result: the stabilising \omterm{def:b2:government-debt:primary-balance}{primary balance} at 3 and at 5 per cent, and the share of a debt-financed tax cut
    that private saving offsets.
\end{enumerate}
\end{problem}
